\section{Introduction}\label{sec:intro}

\subsection{The theorems}\label{ss:intro-theorems}
For a Dirichlet character $\chi$ modulo an integer $N\ge1$ let $L(s,\chi)$ be its $L$-function,
continued meromorphically to $\C$; it is entire unless $\chi$ is principal, in which case it has a
single simple pole, at $s=1$. Let $\zeta(s)$ be the Riemann zeta function. This paper gives a complete
account of a proof of the following two theorems.

\begin{theorem}[The Generalized Riemann Hypothesis for Dirichlet $L$-functions]\label{thm:main-grh}
Let $N\ge1$, let $\chi$ be a Dirichlet character modulo $N$, and let $s\in\C$ with $\re s>1/2$.
Then $L(s,\chi)\neq0$, unless $\chi$ is principal and $s=1$.
\end{theorem}

\begin{theorem}\label{thm:main-rh}
\begin{enumerate}[label=\textup{(\alph*)}]
\item Every zero $s$ of $L(s,\chi)$ with $0<\re s<1$ has $\re s=1/2$, for every Dirichlet character
$\chi$, primitive or not.
\item \textup{(The Riemann Hypothesis.)} Every zero of $\zeta(s)$ other than the trivial zeros
$-2,-4,-6,\dots$ has real part $1/2$.
\end{enumerate}
\end{theorem}

The proof is due to Levent Alp\"oge, working with the AI system Claude \cite{Alpoge2026}; it was
announced together with a complete formalisation in the proof assistant Lean~4
\cite{deMouraUllrich2021}, on top of its mathematical library Mathlib \cite{mathlib2020}. A shorter
formalisation following the same method, the development \texttt{grh-lean} \cite{GRHLean2026}, was
produced at Anthropic in collaboration with Claude and released under the Apache License~2.0; many of
its steps are proved by other arguments than in the original formalisation. This paper is an exposition of the proof as it is carried out in \texttt{grh-lean}. It
introduces no new step: every theorem, proposition and lemma below corresponds to declarations of
that development; \cref{app:lean} lists, section by section, the Lean modules in which they lie.

In \texttt{grh-lean} the five statements are proved in the following form, in the language of
Mathlib. The variable $N$ ranges over all integers $N\ge1$, $\chi$ over all Dirichlet characters
modulo $N$, and $L(s,\chi)$ is Mathlib's $L$-function, the meromorphic continuation of the Dirichlet
series $\sum_{n\ge1}\chi(n)n^{-s}$, regarded as a function on all of $\C$.
\begin{enumerate}[label=\textup{(\arabic*)},leftmargin=2em]
\item For all $s\in\C$: if $\zeta(s)=0$, $s\notin\{-2,-4,-6,\dots\}$ and $s\neq1$, then
$\re s=1/2$. This is Mathlib's definition of the Riemann Hypothesis, proved twice, once under that
name and once in the form stated here.
\item For all $N,\chi,s$: if $\re s>1/2$ and ($\chi\neq1$ or $s\neq1$), then $L(s,\chi)\neq0$.
\item For all $N,\chi,s$: if $0<\re s<1$ and $L(s,\chi)=0$, then $\re s=1/2$.
\item The case of (3) in which $\chi$ is primitive.
\end{enumerate}
\lean{riemannHypothesis, riemannZeta_eq_zero_re, LFunction_ne_zero_of_half_lt_re, LFunction_eq_zero_re, LFunction_eq_zero_re_of_isPrimitive, RiemannZetaZerosOnLine, DirichletLNonvanishingRightOfHalf, DirichletLZerosOnLine, PrimitiveDirichletLZerosOnLine}

Statement~(2) is \cref{thm:main-grh}; statements (1), (3) and (4) together are \cref{thm:main-rh}.
The condition $s\neq1$ in (1) and (2) appears because Mathlib assigns a value to $\zeta$ and to
$L(\,\cdot\,,1)$ at their pole. Statement (3) says nothing about zeros with $\re s\le0$: for a
character that is not primitive, $L(s,\chi)$ differs from the $L$-function of the primitive character
that induces it by finitely many Euler factors $1-\chi_0(p)p^{-s}$, whose zeros lie on $\re s=0$.

\subsection{The Riemann Hypothesis}\label{ss:intro-history}
Riemann's memoir of 1859 \cite{Riemann1859} continued $\zeta(s)=\sum_{n\ge1}n^{-s}$ to the plane,
proved the functional equation relating $s$ and $1-s$, and expressed the number of primes below $x$
by a sum over the zeros of $\zeta$. In the form proved by von Mangoldt \cite{vonMangoldt1895}, for
$x>1$ not a prime power,
\begin{equation}\label{eq:explicit-formula}
  \psi(x)=\sum_{n\le x}\Lam(n)=x-\sum_{\rho}\frac{x^\rho}{\rho}-\log2\pi-\tfrac12\log\bigl(1-x^{-2}\bigr),
\end{equation}
where $\rho$ runs over the zeros of $\zeta$ in the strip $0<\re s<1$, summed symmetrically. Each zero
contributes an oscillation $x^\rho/\rho$ of amplitude $x^{\re\rho}$ and frequency $\im\rho$; these
oscillations are the rhythm of the primes in our title. Riemann wrote that it is ``sehr
wahrscheinlich'' that all roots of his function $\xi(t)$, the completed zeta function on the line
$s=\frac12+it$, are real, that is, that all the zeros $\rho$ have real part $1/2$. If so, then every term has amplitude $\sqrt x$, and the error in
the prime number theorem is as small as the square root of the main term, up to logarithms; and
conversely a bound $\psi(x)=x+O(x^{\beta+\ve})$ for every $\ve>0$ is equivalent to the absence of
zeros in $\re s>\beta$. Hadamard \cite{Hadamard1896} and de la Vall\'ee Poussin
\cite{ValleePoussin1896} proved the prime number theorem in 1896 by showing that $\zeta$ has no
zeros on the line $\re s=1$.

The hypothesis concerns all zeros, while every unconditional result controls a subset of them. Hardy
\cite{Hardy1914} proved that infinitely many zeros lie on the critical line; Selberg
\cite{Selberg1942} that a positive proportion of them do; Levinson \cite{Levinson1974} that more than
one third and Conrey \cite{Conrey1989} that more than two fifths of them do. All zeros with imaginary
part between $0$ and $3\cdot10^{12}$ lie on the line \cite{PlattTrudgian2021}. For the history and the
classical theory see \cite{Edwards1974,Titchmarsh1986}, and for the place of the problem in
mathematics \cite{Bombieri2000,Conrey2003}.

For Dirichlet characters the same picture holds with primes in arithmetic progressions in place of
primes: the zeros of $L(s,\chi)$, $\chi$ modulo $q$, govern the distribution of the primes in the
residue classes modulo $q$, and the Generalized Riemann Hypothesis, the statement of
\cref{thm:main-grh}, gives the error term $O(x^{1/2+\ve})$ in the prime number theorem for every
progression; see \cite{Davenport1980,MontgomeryVaughan2007,IwaniecKowalski2004}.

Over a finite field the analogue of the hypothesis is a theorem. For the zeta function of a curve
over a finite field it was proved by Weil \cite{Weil1948}, who formulated the corresponding
conjectures for varieties of any dimension \cite{Weil1949}; these were proved by Deligne
\cite{Deligne1974,Deligne1980}. Weil's proof rests on a positivity, the Castelnuovo--Severi
inequality on the surface $C\times C$. In 1952 Weil showed that the hypothesis for $\zeta$ itself is
equivalent to the positivity of a quadratic form built from his explicit formula \cite{Weil1952}.
A different, elementary proof was found by Stepanov \cite{Stepanov1969} for hyperelliptic curves
and extended to all curves by Bombieri \cite{Bombieri1974}. It constructs an auxiliary function that vanishes to high order at
the rational points of the curve; the order of vanishing, against the degree of the function, bounds
the number of points, and the bound for all finite extensions gives the hypothesis.

In the monograph \cite{S03} the author analysed the obstructions to transporting these proofs from
function fields to $\Q$, organised around the positivity objects that the known proofs use. Of the two
proofs for curves it took Weil's as its model: in its words, ``it is Weil's proof, not Stepanov's,
that this book takes as the paradigm, because Weil's proof is the one in which the positivity object
is visible'' \cite[p.~279]{S03}; and of Stepanov's method it recorded that ``it produces no
operator, no cohomology and no positive form; and it has resisted every attempt at generalisation to
higher dimension'' \cite[p.~281]{S03}. The proof expounded here takes the other route. It is a
descendant of Stepanov's method, and the positivity it uses is of a different kind: the nonnegativity
of Petersson norms and of a Laplace transform with positive coefficients.

\subsection{The method}\label{ss:intro-method}
We outline the argument; \cref{sec:overview} gives a more detailed plan, and the reading guide at the
end of this section says where each step is proved.

\emph{A one-sided bound suffices.} Fix a prime $d\equiv7\pmod8$, let $\chiK=(\,\cdot\,/d)$ be the
character of $K=\Q(\sqrt{-d})$, and consider for a modulus $q$ the weighted sum
\[
  \sum_{\substack{n\le x\\n\equiv1\ (q)}}\Lam(n)\bigl(1-\chiK(n)\bigr)\Bigl(\frac1n-\frac1x\Bigr)
  =\frac{\log x}{\totient(q)}+c_q+R_q(x).
\]
Its weights are nonnegative, and they vanish on the primes that split in $K$. The Laplace transform of
the remainder $R_q$ has a pole at every zero of an $L$-function modulo $q$, with a residue whose sign
on the real axis is fixed once the zeros of the twisted functions $L(s,\psi\chiK)$ are kept away; a
prime $d$ for which this holds in any given box exists by a zero-density estimate for quadratic twists.
By Landau's theorem on Dirichlet integrals with coefficients of one sign, a bound
$R_q(x)\le x^{\beta-1}$ from \emph{above} only then excludes every zero of every $L(s,\psi)$, $\psi$
modulo $q$, in $\re s>\beta$ (\cref{thm:one-sided-criterion}). It therefore suffices to prove, for one
modulus $q_0$ divisible by any given $N$ and for all primes $d$ in a suitable class,
\begin{equation}\label{eq:intro-one-sided}
  R_{q_0}(x)\le x^{-1/2+\ve}\qquad(x\ \text{large}).
\end{equation}

\emph{A Stepanov argument with theta lifts.} In Stepanov's method an auxiliary function on a curve
vanishes to high order at the rational points, and its degree bounds the number of points. Over $\Z$
there is no differentiation of integers, and the proof replaces the parts of this argument as follows.
\begin{center}
\begin{tabular}{@{}>{\raggedright\arraybackslash}p{.36\textwidth}>{\raggedright\arraybackslash}p{.58\textwidth}@{}}
\toprule
Stepanov's method over $\F_q$ & the proof over $\Q$\\
\midrule
auxiliary function & a chain of $m\asymp\sqrt x$ theta lifts between definite unitary groups\\
high order of vanishing & divisibility, at each prime, of derivatives of theta functions\\
degree bound & archimedean size and height (covolume and index of a lattice)\\
bound for the number of points & the upper bound \eqref{eq:intro-one-sided} for the weighted prime sum\\
\bottomrule
\end{tabular}
\end{center}
The divisibility is where the primes enter. At an inert prime $p$ of the CM field $E$ used below,
the derivatives of order $a$ of theta functions of CM elliptic curves are divisible by a power of
$p$ that is computed by a sum of the shape $\sum_{\nu}\LamK(\nu)\,(y/\nu-\kappa)_+$
(\cref{part:derivatives}); summed over the places of $E$ these divisibilities reproduce the weighted
prime sum, with its main term.

\emph{Three theorems and a main inequality.} For each large $x$ the argument fixes data: scales
$x\le x_1\le b$ with $x_1\ge x^{10^6}$ and $b\ge x_1^{10^6}$, a CM field $E$ with between $x^{11}$ and $x^{12}$
infinite places, ranks of size $x^5$, and three large primes (\cref{sec:data}). It then considers a space $\mathcal T$
of automorphic forms on a definite unitary group, the target, and a family of maps from it, chains of
$m$ theta lifts twisted by anticyclotomic Hecke characters $\xi$. Theorem~A says that a large space of
integral forms has short images, with sizes given by the main amount less the remainder:
\[
  \log\norm{\Theta_1y}^2\le\hmain-m\,\rem+O(fa\log^4b).
\]
Theorem~B says that few translates of these images span almost all of $\mathcal T$, by a second-moment
computation (\cref{part:spanning}). Theorem~C says that integral forms cannot all be short: the
covolume of a spanning family of integral forms, corrected by its index, is at least the main amount
(\cref{part:lower,part:congruences}). It rests on a product formula in which the archimedean sizes are
read through the untwisted chain $\Theta_1$ while the divisibility at each prime $\ell$ is read through
a chain twisted by a character $\xi_\ell$ chosen for that prime, together with a height bound for the
ratio of the two Gram determinants obtained from the Rallis inner product formula and congruences
between theta lifts. Comparing the two bounds, the main amount cancels and
\[
  m\,\rem\le fa\log^5b
\]
(\cref{thm:main-inequality}). Since $m\ge\sqrt x$ and the total remainder $\rem$ is, up to
negligible terms, $a\,\#\Phib\,R_{q_0}(x)$, this is \eqref{eq:intro-one-sided}. The exponent $1/2$
is the reciprocal of the exponent in $m\asymp\sqrt x$: a chain of $m$ lifts plays the part of
vanishing to order $m$.

\subsection{Formalisation and the status of the proof}\label{ss:intro-status}
The development \texttt{grh-lean} \cite{GRHLean2026} consists of a readable edition of about $300{,}000$
lines of Lean code and a compact edition of about $76{,}000$ lines, which prove the same five
statements. Its authors report that every file compiles with Lean v4.35.0 and Mathlib at commit
\texttt{56a00585}, without \texttt{sorry}; that the five theorems depend only on the three axioms
\texttt{propext}, \texttt{Classical.choice} and \texttt{Quot.sound}; and that an export of the five
theorems with everything they depend on is accepted by three independent proof checkers. They state
that no mathematician has read the proof through and that no paper accompanies it.

We rebuilt the readable edition at the pinned versions and checked the axioms and the statements
ourselves; \cref{app:verification} records how, with the hashes of the sources. Our own reading of the
Lean text found the argument presented below. Its mathematical correctness is what the formal
verification establishes; the present exposition has not been refereed, and where it and the Lean text
differ the Lean text is authoritative and the difference is an error of this paper.

The development is organised in two parts. One part, about $177{,}000$ lines, proves results from the
literature that Mathlib lacks, often in a weaker form that suffices here and with its own definitions:
class field theory, $\ell$-adic and $p$-adic cohomology, Galois representations attached to automorphic
forms, the Siegel--Weil and Rallis formulas, and others. The other part, about $122{,}000$ lines, carries
out the argument. This paper proves the second part in full. For the first, \cref{app:background} states
precisely the form in which each result is proved and used, with its sources; we do not reproduce
those proofs.

\subsection{Organisation}\label{ss:intro-organisation}
\Cref{sec:overview} sets out the chain of the proof and the place of each step in the paper
(\cref{fig:chain,tab:parts}), and explains the argument at the level of the main objects and their sizes, and
\cref{sec:notation} fixes the notation. \Cref{part:reduction} reduces the theorems to the three theorems
A, B and C and proves the main inequality. \Cref{part:derivatives} computes the divisibility of
derivatives of theta functions. \Cref{part:upper} constructs the composite lift and proves Theorem~A.
\Cref{part:spanning} proves Theorem~B. \Cref{part:lower,part:congruences} prove Theorem~C: the first
treats determinants of chains of theta lifts and the cost of a single lift, the second the congruences
and Selmer groups that bound the denominators. \Cref{app:ds} contains an estimate for discrete series
coefficients used in \cref{part:spanning,part:lower}; \cref{app:background} the background results;
\cref{app:lean} the correspondence with the Lean declarations; \cref{app:constants} the constants and
exponents; and \cref{app:verification} the record of the formal verification.

The names of the Lean declarations that state and prove each statement are listed, statement by
statement, in the supplement to the paper (\cref{app:lean}). A reader who wishes to check a step can
therefore read the corresponding Lean statement, which is complete and precise, next to the proof
given here, which explains it.
